\documentclass[10pt,a4paper]{amsart}
\usepackage{amsmath,amssymb,amsthm,graphicx}
\usepackage{verbatim}
\newenvironment{reference}{\comment}{\endcomment}
%\newenvironment{reference}{}{}
\newenvironment{slogan}{\comment}{\endcomment}
\newenvironment{history}{\comment}{\endcomment}

% For commutative diagrams we use Xy-pic
\usepackage[all]{xy}

% We use 2cell for 2-commutative diagrams.
\xyoption{2cell}
% Theorem environments.
%
\theoremstyle{plain}
\newtheorem{theorem}[subsection]{Theorem}
\newtheorem{proposition}[subsection]{Proposition}
\newtheorem{lemma}[subsection]{Lemma}

\theoremstyle{definition}
\newtheorem{definition}[subsection]{Definition}
\newtheorem{example}[subsection]{Example}
\newtheorem{exercise}[subsection]{Exercise}
\newtheorem{situation}[subsection]{Situation}

\theoremstyle{remark}
\newtheorem{remark}[subsection]{Remark}
\newtheorem{remarks}[subsection]{Remarks}

\numberwithin{equation}{subsection}

% Macros
%
\def\lim{\mathop{\mathrm{lim}}\nolimits}
\def\colim{\mathop{\mathrm{colim}}\nolimits}
\def\Spec{\mathop{\mathrm{Spec}}}
\def\Hom{\mathop{\mathrm{Hom}}\nolimits}
\def\Ext{\mathop{\mathrm{Ext}}\nolimits}
\def\SheafHom{\mathop{\mathcal{H}\!\mathit{om}}\nolimits}
\def\SheafExt{\mathop{\mathcal{E}\!\mathit{xt}}\nolimits}
\def\Sch{\mathit{Sch}}
\def\Mor{\mathop{\mathrm{Mor}}\nolimits}
\def\Ob{\mathop{\mathrm{Ob}}\nolimits}
\def\Sh{\mathop{\mathit{Sh}}\nolimits}
\def\NL{\mathop{N\!L}\nolimits}
\def\CH{\mathop{\mathrm{CH}}\nolimits}
\def\proetale{{pro\text{-}\acute{e}tale}}
\def\etale{{\acute{e}tale}}
\def\QCoh{\mathit{QCoh}}
\def\Ker{\mathop{\mathrm{Ker}}}
\def\Im{\mathop{\mathrm{Im}}}
\def\Coker{\mathop{\mathrm{Coker}}}
\def\Coim{\mathop{\mathrm{Coim}}}

% Boxtimes
%
\DeclareMathSymbol{\boxtimes}{\mathbin}{AMSa}{"02}

%
% Macros for moduli stacks/spaces
%
\def\QCohstack{\mathcal{QC}\!\mathit{oh}}
\def\Cohstack{\mathcal{C}\!\mathit{oh}}
\def\Spacesstack{\mathcal{S}\!\mathit{paces}}
\def\Quotfunctor{\mathrm{Quot}}
\def\Hilbfunctor{\mathrm{Hilb}}
\def\Curvesstack{\mathcal{C}\!\mathit{urves}}
\def\Polarizedstack{\mathcal{P}\!\mathit{olarized}}
\def\Complexesstack{\mathcal{C}\!\mathit{omplexes}}
% \Pic is the operator that assigns to X its picard group, usage \Pic(X)
% \Picardstack_{X/B} denotes the Picard stack of X over B
% \Picardfunctor_{X/B} denotes the Picard functor of X over B
\def\Pic{\mathop{\mathrm{Pic}}\nolimits}
\def\Picardstack{\mathcal{P}\!\mathit{ic}}
\def\Picardfunctor{\mathrm{Pic}}
\def\Deformationcategory{\mathcal{D}\!\mathit{ef}}
\usepackage{cleveref}
\crefname{originallemma}{original lemma}{original lemmas}
\Crefname{originallemma}{Original lemma}{Original lemmas}
\makeatletter
\newlabel{divisors-lemma-blow-up-effective-Cartier-divisor}{{0807}{source}{}{sourcepointer-1}{}}
\newlabel{divisors-lemma-blow-up-effective-Cartier-divisor@cref}{{[originallemma][0][]0807}{[1][0][]source}{}{}{}}
\makeatother
\crefname{originallemma}{original lemma}{original lemmas}
\Crefname{originallemma}{Original lemma}{Original lemmas}
\makeatletter
\newlabel{divisors-lemma-strict-transform}{{080E}{source}{}{sourcepointer-2}{}}
\newlabel{divisors-lemma-strict-transform@cref}{{[originallemma][0][]080E}{[1][0][]source}{}{}{}}
\makeatother
\title{Behaviour of the strict transforms of two curves when one of them is regular at the blown up point}
\author{The Stacks Project Authors}
\date{stacks-project:resolve.tex:local-snapshot-20261012; GNU FDL 1.2 or later}
\begin{document}
\maketitle
\noindent\textit{Editorial note.} All delivered mathematical text is the author's own, copied from the pinned Stacks Project chapter file resolve.tex; no statement or proof step was rewritten, reordered or removed.\par\medskip
\noindent\textit{Editorial note.} The counted proof is the author's complete argument: the choice of a local equation of the regular curve bringing the maximal ideal into a convenient shape, the identification of the strict transform through the blowup of an effective Cartier divisor and the explicit map from the blowup algebra, and the computation of the two intersection multiplicities.\par\medskip
\noindent\textit{Editorial note.} The label lemma-blowup-nonsingular-curves-meeting-at-point names the counted statement, and the proof environment immediately following it in the source is delivered verbatim as its proof.\par\medskip
\noindent\textit{Editorial note.} The counted statement is carried by the author's own lemma environment in the chapter source.\par\medskip
\noindent\textit{Editorial note.} Pre-open gate: python operations/agent-authoring/tools/precard\_gate.py --source resolve.tex --label lemma-blowup-nonsingular-curves-meeting-at-point --start 4345 --end 4392 --local-id 28096 returned RESULT=OK (L28096 free, occupancy FREE, public interval clear). Pre-assembly tag self-check: canonical tag resolved and not present in reference/public-items.jsonl.\par\medskip
\noindent\textit{Editorial note.} The statement and proof are copied from the pinned chapter file; the cited inputs are kept as external pointers with their canonical Stacks tags.\par\medskip
\noindent\textit{Allowed standard prerequisites (editorial).}\begin{itemize}
\item Strict transform as the blowup of the effective Cartier divisor. Divisors, the lemma identifying the strict transform of a divisor under a blowup, and the lemma that blowing up an effective Cartier divisor does not change the scheme. Supplies the identification of the strict transform of the regular curve with the blowup of that curve in the point, hence the isomorphism statement.
\item Description of the blowup algebra and the ideal of the exceptional divisor. Standard description of the affine blowup algebra generated by the quotients of the maximal ideal by a chosen element, with the image of that element cutting out the exceptional divisor. Supplies the algebra on which the comparison map is defined and the computation of the multiplicity against the exceptional fibre.
\end{itemize}\medskip
\ifdefined\hypertarget\hypertarget{sourcepointer-1}{}\fi
\noindent\textit{Source reference.} Original lemma 0807: sources/stacks-chow/tags.txt, source lines 10172--10172. Prerequisite pointer invoked at proof line(s) 4372; it names a standard fact of a sibling chapter that the counted argument uses, and every step of the counted proof built on it is retained verbatim in the delivered body.\par\medskip
\ifdefined\hypertarget\hypertarget{sourcepointer-2}{}\fi
\noindent\textit{Source reference.} Original lemma 080E: sources/stacks-chow/tags.txt, source lines 10179--10179. Prerequisite pointer invoked at proof line(s) 4370; it names a standard fact of a sibling chapter that the counted argument uses, and every step of the counted proof built on it is retained verbatim in the delivered body.\par\medskip
\begin{lemma}
\label{lemma-blowup-nonsingular-curves-meeting-at-point}
In the situation above let $X' \to X$ be the blowing up of $X$ in $p$.
Let $Y', Z' \subset X'$ be the strict transforms of $Y, Z$.
If $\mathcal{O}_{Y, p}$ is regular, then
\begin{enumerate}
\item $Y' \to Y$ is an isomorphism,
\item $Y'$ meets the exceptional fibre $E \subset X'$ in one point
$q$ and $m_q(Y \cap E) = 1$,
\item if $q \in Z'$ too, then $m_q(Y \cap Z') < m_p(Y \cap Z)$.
\end{enumerate}
\end{lemma}

\begin{proof}
Since $\mathcal{O}_{X, p} \to \mathcal{O}_{Y, p}$ is surjective and
$\mathcal{O}_{Y, p}$ is a discrete valuation ring, we can pick an
element $x_1 \in \mathfrak m_p$ mapping to a uniformizer in
$\mathcal{O}_{Y, p}$. Choose an affine open $U = \Spec(A)$ containing
$p$ such that $x_1 \in A$. Let $\mathfrak m \subset A$ be the
maximal ideal corresponding to $p$. Let $I, J \subset A$
be the ideals defining $Y, Z$ in $\Spec(A)$. After shrinking $U$
we may assume that $\mathfrak m = I + (x_1)$, in other words,
that $V(x_1) \cap U \cap Y = \{p\}$ scheme theoretically.
We conclude that $p$ is an effective Cartier divisor on $Y$ and
since $Y'$ is the blowing up of $Y$ in $p$
(Divisors, Lemma \ref{divisors-lemma-strict-transform})
we see that $Y' \to Y$ is an isomorphism by
Divisors, Lemma \ref{divisors-lemma-blow-up-effective-Cartier-divisor}.
The relationship $\mathfrak m = I + (x_1)$ implies that
$\mathfrak m^n \subset I + (x_1^n)$ hence we can define a map
$$
\psi : A[\textstyle{\frac{\mathfrak m}{x_1}}] \longrightarrow A/I
$$
by sending $y/x_1^n \in A[\frac{\mathfrak m}{x_1}]$
to the class of $a$ in $A/I$ where $a$
is chosen such that $y \equiv ax_1^n \bmod I$.
Then $\psi$ corresponds to the morphism of $Y \cap U$ into $X'$
over $U$ given by $Y' \cong Y$. Since the image of $x_1$
in $A[\frac{\mathfrak m}{x_1}]$ cuts out the exceptional divisor
we conclude that $m_q(Y', E) = 1$. Finally, since
$J \subset \mathfrak m$ implies that the ideal
$J' \subset A[\frac{\mathfrak m}{x_1}]$ certainly
contains the elements $f/x_1$ for $f \in J$.
Thus if we choose $f \in J$ whose image $\overline{f}$ in $A/I$ has
minimal valuation equal to $m_p(Y \cap Z)$, then we see that
$\psi(f/x_1) = \overline{f}/x_1$ in $A/I$ has valuation one less
proving the last part of the lemma.
\end{proof}

\end{document}
